% The bibliography is embedded below so this single file compiles on its own.
% LaTeX writes it to references.bib at compile time; BibTeX then builds the reference list.
\begin{filecontents*}[overwrite]{references.bib}
% Entries marked "verified" were checked against the primary page (arXiv, publisher or owner site)
% during this project. Classic references without that mark were entered from standard bibliographic
% knowledge; please confirm page numbers and DOIs against the publisher before submission.

@misc{srinivasan2026snap2,
  author = {Srinath Srinivasan and Tim Menzies},
  title = {Better Together, in the Right Order: Classical-then-{LLM} Optimization for {SE}},
  year = {2026},
  eprint = {2607.02583},
  archivePrefix = {arXiv},
  url = {https://arxiv.org/abs/2607.02583},
  verification = {checked against primary page}
}

@misc{menzies2026ezr,
  author = {Tim Menzies and Srinath Srinivasan and Kishan Ganguly},
  title = {Can {AI} be Easy? {L}essons Learned from the {EZR.py} Toolkit},
  year = {2026},
  eprint = {2606.03640},
  archivePrefix = {arXiv},
  url = {https://arxiv.org/abs/2606.03640},
  note = {Toolkit version 0.9.4},
  verification = {checked against arXiv}
}

@misc{menzies2025moot,
  author = {Tim Menzies and Tao Chen and Yulong Ye and Kishan Kumar Ganguly and Amirali Rayegan and Srinath Srinivasan and Andre Lustosa},
  title = {{MOOT}: a Repository of Many Multi-Objective Optimization Tasks},
  year = {2025},
  eprint = {2511.16882},
  archivePrefix = {arXiv},
  url = {https://arxiv.org/abs/2511.16882},
  verification = {checked against primary page}
}

@misc{ganguly2026tournament,
  author = {Kishan Kumar Ganguly and Tim Menzies},
  title = {Which Optimizer, At What Budget? {A} Tournament of Optimizers for Search-Based {SE}},
  year = {2026},
  eprint = {2607.11705},
  archivePrefix = {arXiv},
  url = {https://arxiv.org/abs/2607.11705},
  verification = {checked against primary page}
}

@misc{rayegan2026instability,
  author = {Amirali Rayegan and Lunxiao Li and Tim Menzies},
  title = {Is Model Instability just Noise to be Tolerated or a Property that can be Managed?},
  year = {2026},
  eprint = {2607.10420},
  archivePrefix = {arXiv},
  url = {https://arxiv.org/abs/2607.10420},
  verification = {checked against primary page}
}

@inproceedings{liu2024llambo,
  author = {Tennison Liu and Nicol{\'a}s Astorga and Nabeel Seedat and Mihaela van der Schaar},
  title = {Large Language Models to Enhance {B}ayesian Optimization},
  booktitle = {International Conference on Learning Representations (ICLR)},
  year = {2024},
  url = {https://arxiv.org/abs/2402.03921},
  verification = {checked against primary page}
}

@misc{rychert2025reproduction,
  author = {Adam Rychert and Gasper Spagnolo and Evgenii Posashkov},
  title = {Reproducibility Study of Large Language Model {B}ayesian Optimization},
  year = {2025},
  eprint = {2511.18891},
  archivePrefix = {arXiv},
  url = {https://arxiv.org/abs/2511.18891},
  verification = {checked against primary page}
}

@misc{rodrigues2026worth,
  author = {Carson Rodrigues and Oysturn Vas and Isaiah Abner DCosta and Nithish Kumar Prabhakaran},
  title = {When Is an {LLM} Worth It for Hyperparameter Optimization? {A} Budget-Matched Study on Tabular Data Finds the Warm-Start Is a Default Configuration, Not the Model},
  year = {2026},
  eprint = {2606.21641},
  archivePrefix = {arXiv},
  url = {https://arxiv.org/abs/2606.21641},
  verification = {checked against primary page}
}

@misc{cisse2025bora,
  author = {Abdoulatif Ciss{\'e} and Xenophon Evangelopoulos and Vladimir V. Gusev and Andrew I. Cooper},
  title = {Language-Based {B}ayesian Optimization Research Assistant ({BORA})},
  year = {2025},
  eprint = {2501.16224},
  archivePrefix = {arXiv},
  url = {https://arxiv.org/abs/2501.16224},
  verification = {checked against primary page}
}

@misc{xu2026lbmcts,
  author = {Beicheng Xu and Weitong Qian and Lingching Tung and Yupeng Lu and Bin Cui},
  title = {Tree-Structured Synergy of Large Language Models and {B}ayesian Optimization for Efficient {CASH}},
  year = {2026},
  eprint = {2601.12355},
  archivePrefix = {arXiv},
  url = {https://arxiv.org/abs/2601.12355},
  verification = {checked against primary page}
}

@misc{ngo2026lmabo,
  author = {Giang Ngo and Dat Phan Trong and Dang Nguyen and Sunil Gupta and Svetha Venkatesh},
  title = {Adaptive Acquisition Selection for {B}ayesian Optimization with Large Language Models},
  year = {2026},
  eprint = {2602.07904},
  archivePrefix = {arXiv},
  url = {https://arxiv.org/abs/2602.07904},
  verification = {checked against primary page}
}

@misc{chen2023frugalgpt,
  author = {Lingjiao Chen and Matei Zaharia and James Zou},
  title = {{FrugalGPT}: How to Use Large Language Models While Reducing Cost and Improving Performance},
  year = {2023},
  eprint = {2305.05176},
  archivePrefix = {arXiv},
  url = {https://arxiv.org/abs/2305.05176}
}

@inproceedings{aggarwal2023automix,
  author = {Pranjal Aggarwal and Aman Madaan and Ankit Anand and Srividya Pranavi Potharaju and Swaroop Mishra and Pei Zhou and Aditya Gupta and Dheeraj Rajagopal and Karthik Kappaganthu and Yiming Yang and Shyam Upadhyay and Manaal Faruqui and Mausam},
  title = {{AutoMix}: Automatically Mixing Language Models},
  booktitle = {Advances in Neural Information Processing Systems (NeurIPS)},
  year = {2024},
  url = {https://arxiv.org/abs/2310.12963},
  verification = {checked against primary page}
}

@misc{ong2024routellm,
  author = {Isaac Ong and Amjad Almahairi and Vincent Wu and Wei-Lin Chiang and Tianhao Wu and Joseph E. Gonzalez and M. Waleed Kadous and Ion Stoica},
  title = {{RouteLLM}: Learning to Route {LLMs} with Preference Data},
  year = {2024},
  eprint = {2406.18665},
  archivePrefix = {arXiv},
  url = {https://arxiv.org/abs/2406.18665},
  verification = {checked against primary page}
}

@misc{huang2026scope,
  author = {Yiqian Huang and Shiqi Zhang and Tianyuan Jin and Xiaokui Xiao},
  title = {{SCOPE}: Cost-Efficient Model Selection for Compound {AI} Systems under Quality Constraints},
  year = {2026},
  eprint = {2606.00774},
  archivePrefix = {arXiv},
  url = {https://arxiv.org/abs/2606.00774},
  verification = {checked against primary page}
}

@inproceedings{mozannar2020defer,
  author = {Hussein Mozannar and David Sontag},
  title = {Consistent Estimators for Learning to Defer to an Expert},
  booktitle = {Proceedings of the 37th International Conference on Machine Learning (ICML)},
  year = {2020},
  url = {https://arxiv.org/abs/2006.01862},
  verification = {checked against primary page}
}

@inproceedings{geifman2017selective,
  author = {Yonatan Geifman and Ran El-Yaniv},
  title = {Selective Classification for Deep Neural Networks},
  booktitle = {Advances in Neural Information Processing Systems (NeurIPS)},
  year = {2017},
  url = {https://arxiv.org/abs/1705.08500}
}

@article{harman2001sbse,
  author = {Mark Harman and Bryan F. Jones},
  title = {Search-Based Software Engineering},
  journal = {Information and Software Technology},
  volume = {43},
  number = {14},
  pages = {833--839},
  year = {2001}
}

@article{nair2020flash,
  author = {Vivek Nair and Zhe Yu and Tim Menzies and Norbert Siegmund and Sven Apel},
  title = {Finding Faster Configurations Using {FLASH}},
  journal = {IEEE Transactions on Software Engineering},
  volume = {46},
  number = {7},
  year = {2020}
}

@inproceedings{siegmund2015influence,
  author = {Norbert Siegmund and Alexander Grebhahn and Sven Apel and Christian K{\"a}stner},
  title = {Performance-Influence Models for Highly Configurable Systems},
  booktitle = {Proceedings of the 10th Joint Meeting on Foundations of Software Engineering (ESEC/FSE)},
  pages = {284--294},
  year = {2015}
}

@inproceedings{ha2019deepperf,
  author = {Huong Ha and Hongyu Zhang},
  title = {{DeepPerf}: Performance Prediction for Configurable Software with Deep Sparse Neural Network},
  booktitle = {Proceedings of the 41st International Conference on Software Engineering (ICSE)},
  year = {2019}
}

@article{kaltenecker2023evolution,
  author = {Christian Kaltenecker and Stefan M{\"u}hlbauer and Alexander Grebhahn and Norbert Siegmund and Sven Apel},
  title = {Performance Evolution of Configurable Software Systems: An Empirical Study},
  journal = {Empirical Software Engineering},
  year = {2023},
  doi = {10.1007/s10664-023-10338-3},
  verification = {checked against primary page}
}

@inproceedings{fekry2020tuneful,
  author = {Ayat Fekry and Lucian Carata and Thomas Pasquier and Andrew Rice and Andy Hopper},
  title = {To Tune or Not to Tune? {I}n Search of Optimal Configurations for Data Analytics},
  booktitle = {Proceedings of the 26th ACM SIGKDD Conference on Knowledge Discovery and Data Mining (KDD)},
  year = {2020},
  doi = {10.1145/3394486.3403299},
  verification = {checked against primary page}
}

@misc{hsu2018scout,
  author = {Chin-Jung Hsu and Vivek Nair and Tim Menzies and Vincent W. Freeh},
  title = {Scout: An Experienced Guide to Find the Best Cloud Configuration},
  year = {2018},
  eprint = {1803.01296},
  archivePrefix = {arXiv},
  url = {https://arxiv.org/abs/1803.01296},
  verification = {checked against primary page}
}

@article{jones1998ego,
  author = {Donald R. Jones and Matthias Schonlau and William J. Welch},
  title = {Efficient Global Optimization of Expensive Black-Box Functions},
  journal = {Journal of Global Optimization},
  volume = {13},
  number = {4},
  pages = {455--492},
  year = {1998}
}

@inproceedings{barbosa2022cvc5,
  author = {Haniel Barbosa and others},
  title = {cvc5: A Versatile and Industrial-Strength {SMT} Solver},
  booktitle = {Tools and Algorithms for the Construction and Analysis of Systems (TACAS)},
  year = {2022}
}

@misc{perron2023ortools,
  author = {Laurent Perron and Vincent Furnon},
  title = {{OR-Tools}},
  howpublished = {Google, \url{https://developers.google.com/optimization}},
  year = {2025},
  note = {Version 9.14}
}

@misc{gallant2016ripgrep,
  author = {Andrew Gallant},
  title = {ripgrep},
  howpublished = {\url{https://github.com/BurntSushi/ripgrep}},
  year = {2016}
}

@article{malkov2020hnsw,
  author = {Yu. A. Malkov and D. A. Yashunin},
  title = {Efficient and Robust Approximate Nearest Neighbor Search Using Hierarchical Navigable Small World Graphs},
  journal = {IEEE Transactions on Pattern Analysis and Machine Intelligence},
  volume = {42},
  number = {4},
  pages = {824--836},
  year = {2020}
}

@misc{vink2020polars,
  author = {Ritchie Vink and others},
  title = {Polars},
  howpublished = {\url{https://pola.rs}},
  year = {2025},
  note = {Version 1.35.2}
}

@inproceedings{chen2016xgboost,
  author = {Tianqi Chen and Carlos Guestrin},
  title = {{XGBoost}: A Scalable Tree Boosting System},
  booktitle = {Proceedings of the 22nd ACM SIGKDD International Conference on Knowledge Discovery and Data Mining (KDD)},
  pages = {785--794},
  year = {2016}
}

@misc{gerganov2023llamacpp,
  author = {Georgi Gerganov and others},
  title = {llama.cpp},
  howpublished = {\url{https://github.com/ggml-org/llama.cpp}},
  year = {2025},
  note = {Build b11146}
}

@misc{huggingface2025smollm3,
  author = {Elie Bakouch and others},
  title = {{SmolLM3}: smol, multilingual, long-context reasoner},
  howpublished = {Hugging Face blog, \url{https://huggingface.co/blog/smollm3}},
  year = {2025},
  note = {Published 8 July 2025},
  verification = {checked against source page}
}

@misc{qwen2025qwen3,
  author = {An Yang and others},
  title = {{Qwen3} Technical Report},
  year = {2025},
  eprint = {2505.09388},
  archivePrefix = {arXiv},
  url = {https://arxiv.org/abs/2505.09388},
  verification = {checked against primary page}
}

@misc{openai2025gptoss,
  author = {{OpenAI}},
  title = {gpt-oss-120b \& gpt-oss-20b Model Card},
  year = {2025},
  eprint = {2508.10925},
  archivePrefix = {arXiv},
  url = {https://arxiv.org/abs/2508.10925},
  verification = {checked against primary page}
}

@misc{openrouter2025,
  author = {{OpenRouter}},
  title = {{OpenRouter}: A Unified Interface for {LLMs}},
  howpublished = {\url{https://openrouter.ai}},
  year = {2025}
}

@article{gelman2014forking,
  author = {Andrew Gelman and Eric Loken},
  title = {The Statistical Crisis in Science},
  journal = {American Scientist},
  volume = {102},
  number = {6},
  pages = {460--465},
  year = {2014}
}
\end{filecontents*}

\documentclass[11pt]{article}

\usepackage[utf8]{inputenc}
\usepackage[T1]{fontenc}
\usepackage{lmodern}
\usepackage[margin=1in]{geometry}
\usepackage{amsmath}
\usepackage{booktabs}
\usepackage{multirow}
\usepackage{array}
\usepackage{xcolor}
\usepackage[numbers,sort&compress]{natbib}
\usepackage[hidelinks]{hyperref}
\usepackage{microtype}

\newcommand{\RQ}[1]{\textbf{RQ#1}}

\title{When Is an LLM Worth Calling?\\[4pt]
\large Rare, Concentrated Gains from Escalating Classical Configuration Search to Large Language Models}

\author{Mayank Dhingra\\
\small [Affiliation]\\
\small [Email]}

\date{}

\begin{document}
\maketitle

\begin{abstract}
Recent work in search-based software engineering suggests that a classical optimizer and a large language model (LLM) work well in sequence: the classical method spends the first part of the evaluation budget and the LLM proposes configurations for the rest. This raises a practical question that has not been measured carefully. After the classical stage, is the LLM worth calling at all, and can a cheap controller tell in advance? We studied this question on 70 configuration-optimization cases drawn from eight software engines (seven ecosystems), and on six natively executed systems. From an identical ten-evaluation prefix, we compared LLM continuations with several cheap classical continuations under a charged budget of 20 evaluations. We tested three local open-weight models (3B, 8B and 14B parameters) and, in a pre-registered final experiment, the model used by the original classical-then-LLM study, gpt-oss-120b, both one-shot and inside an iterative loop modelled on that study. Every model lost to continued classical search on average, by 3.1\% to 5.2\%. Most apparent LLM wins were not specific to the LLM: random search, a neighbourhood rule or Gaussian-process expected improvement usually did as well. The 3B and 8B models produced no case in which the LLM beat every cheap alternative, and the 14B model produced two. gpt-oss-120b produced a few such cases, spread over one or two of the seven ecosystems depending on the run, and one case improved in all three of its runs (by up to 31\% under the iterative loop). Because these wins are rare and cluster in one ecosystem, none of the controllers we tried (learned benefit predictors, uncertainty thresholds, and rules adapted from BORA and LB-MCTS) did better than never calling the LLM. We argue that the benefit of escalation should be measured against the best cheap alternative rather than against the incumbent optimizer, and that routing studies need to show benefit in several independent system groups before a router can be evaluated at all.
\end{abstract}

\noindent\textbf{Keywords:} search-based software engineering, configuration optimization, large language models, active learning, model routing, negative results

\section{Introduction}

Finding a good configuration for a configurable software system is expensive because each evaluation means building, running or benchmarking the system. Work in search-based software engineering has therefore focused on optimizers that do well with very small labeling budgets, often a few dozen evaluations or fewer \cite{harman2001sbse,nair2020flash,menzies2026ezr}. Benchmark collections such as MOOT make it possible to compare such optimizers across many tasks \cite{menzies2025moot}, and tournament-style comparisons show that which optimizer wins depends heavily on the budget \cite{ganguly2026tournament}.

Large language models have recently been added to this toolbox. LLAMBO showed that an LLM can warm-start, model and propose candidates inside Bayesian optimization \cite{liu2024llambo}, although a later reproduction reported its benefits with a 70B-parameter model and ran into problems with smaller ones \cite{rychert2025reproduction}. Closer to software engineering, Srinivasan and Menzies proposed SNAP2, which runs a classical active learner for the first half of a 20-label budget and hands the second half to an LLM \cite{srinivasan2026snap2}. Their results favour the classical-then-LLM order, and they name \emph{conditional} escalation, calling the LLM only when it is likely to help, as future work.

That suggestion is the starting point for this paper. If the LLM stage helps only sometimes, a controller that looks at the classical prefix and predicts whether escalation will pay off could save most LLM calls without losing much. Routing and deferral methods for LLMs make this plausible \cite{chen2023frugalgpt,ong2024routellm,mozannar2020defer}, and several Bayesian-optimization systems already decide when to involve an LLM from signals such as uncertainty, plateaus or surrogate reliability \cite{cisse2025bora,xu2026lbmcts,ngo2026lmabo}.

We set out to build such a controller and did not succeed. The reasons turned out to be more informative than the controller would have been. We compared LLM continuations with cheap classical continuations from the same saved prefix, under the same charged budget, on 70 recorded-data cases and six natively executed systems. We tested three local models and, in a final pre-registered experiment, the same hosted model SNAP2 used, both in our one-shot interface and in an iterative loop modelled on SNAP2's. We found that:

\begin{itemize}
  \item Calling the LLM after the classical prefix loses on average against continuing the classical search, for every model we tried (Section~\ref{sec:rq1}).
  \item Most LLM wins are not due to the LLM. When an LLM continuation beats the incumbent optimizer, random search or GP-based expected improvement from the same prefix usually does as well. Measured against the best cheap alternative, the 3B and 8B models never produced an LLM-specific win, and the 14B model produced two, both on Spark (Section~\ref{sec:rq2}).
  \item Stronger models help a little. gpt-oss-120b produced a handful of LLM-specific wins, one of which appeared in all three of its runs and was largest under the iterative loop (Section~\ref{sec:rq3}).
  \item The wins that exist are rare and concentrated in a single ecosystem. Under leave-one-group-out evaluation, a controller then either has no positive examples to learn from or no positive cases to find. Every controller we tried either learned never to call the LLM, which was the correct decision rather than a failure of the learner, or called it and lost (Section~\ref{sec:rq4}).
\end{itemize}

The paper makes three contributions. First, it reports a paired, budget-charged comparison of LLM and classical continuations across two cohorts and five model configurations, with all failures and fallbacks kept in the denominators. Second, it separates \emph{LLM-specific} benefit from the benefit of simply switching away from the incumbent optimizer, and shows that the distinction changes the conclusions. Third, it explains why escalation routers are hard to evaluate on current benchmarks: benefit that is rare and concentrated in one group makes cross-system evaluation uninformative, so zero-call behaviour should be read as correct abstention.

We should be clear about how the study was carried out. It went through many iterations on the same exposed systems (Section~\ref{sec:history}). Only the last two experiments were frozen in full before data collection, with decision rules fixed in advance. We report the earlier experiments because they carry most of the controller evidence, but we treat them as exploratory.

\section{Background and Related Work}

\subsection{Configuration optimization on small budgets}

Performance-influence models and their successors predict how configuration options affect runtime or throughput \cite{siegmund2015influence,ha2019deepperf}. Optimizers built on such models, or on simpler nearest-neighbour and centroid heuristics, try to find a good configuration while labeling as few as possible \cite{nair2020flash,menzies2026ezr}. The EZR toolkit is a recent example that favours simple acquisition functions and very small budgets \cite{menzies2026ezr}. Two lessons from this literature matter here. First, cheap optimizers are strong baselines at small budgets \cite{ganguly2026tournament}. Second, repeated runs of the same method vary, so single-run comparisons can mislead \cite{rayegan2026instability}.

\subsection{LLMs inside the optimization loop}

LLAMBO uses an LLM for warm-starting, surrogate modelling and candidate sampling in Bayesian optimization \cite{liu2024llambo}. A reproduction study reports benefits with a 70B-parameter model but weaker surrogate regression and problems with smaller backbones \cite{rychert2025reproduction}. In tabular hyperparameter optimization, a budget-matched study found that much of the LLM's apparent advantage came from its warm start being close to a good default configuration rather than from the model's reasoning \cite{rodrigues2026worth}. Our attribution analysis reaches a similar conclusion in a different setting.

Several methods decide \emph{when} to use the LLM. BORA triggers LLM involvement when the Gaussian process is uncertain or progress has plateaued, and adjusts its trust in the LLM from past interventions \cite{cisse2025bora}. LB-MCTS sets the probability of an LLM proposal from the cross-validated rank reliability of the surrogate \cite{xu2026lbmcts}. LMABO lets an LLM choose the acquisition function at each step from the optimization state \cite{ngo2026lmabo}. SNAP2 fixes the hand-off point at half the budget and asks for two proposals per call, with feedback on the measured trajectory and on collisions \cite{srinivasan2026snap2}. None of these studies evaluates a pre-decision controller for a single escalation against the best cheap alternative on held-out system groups, which is the question we ask.

\subsection{Routing and deferral}

FrugalGPT cascades queries through progressively more expensive models \cite{chen2023frugalgpt}. AutoMix decides when to escalate from a small to a large model \cite{aggarwal2023automix}. RouteLLM learns a router from preference data \cite{ong2024routellm}, and SCOPE selects models for compound systems under quality constraints \cite{huang2026scope}. Learning to defer and selective classification study when a predictor should hand a case to an expert or abstain \cite{mozannar2020defer,geifman2017selective}. These methods route individual queries whose answers can be scored directly. Our decision is different. The value of escalation depends on a counterfactual search trajectory under a charged evaluation budget, and it has to be compared with other ways of spending the same budget. This difference turns out to matter.

\section{Study Design}

\subsection{Research questions}

\begin{description}
  \item[\RQ{1} (Incremental value).] From the same ten-evaluation prefix, does an LLM continuation find better configurations than continuing the classical search?
  \item[\RQ{2} (Attribution).] When an LLM continuation wins, would a cheap classical alternative (random search, a neighbourhood rule, or GP-based expected improvement) have done as well?
  \item[\RQ{3} (Model and interaction).] Do a larger local model, and SNAP2's own model used one-shot or iteratively, change the answers to \RQ{1} and \RQ{2}?
  \item[\RQ{4} (Routing).] Can a controller that sees only the prefix predict which escalations will pay off on systems it was not trained on?
\end{description}

\subsection{Tasks}

\paragraph{Recorded cohort.} Our main cohort contains 70 cases built from published per-configuration measurements (Table~\ref{tab:cohort}).
\begin{itemize}
  \item Five families come from the data released with DeepPerf \cite{ha2019deepperf}, which in turn derive from SPLConqueror measurements \cite{siegmund2015influence}: BerkeleyDB (C edition), DUNE, HIPAcc, LLVM and SaC.
  \item OpenVPN comes from the Performance Evolution study \cite{kaltenecker2023evolution}; its objective is throughput, so higher is better.
  \item Five Spark workloads come from the Tuneful data \cite{fekry2020tuneful}, and three Hadoop MapReduce workloads come from the Scout data \cite{hsu2018scout}. For Hadoop, runs that did not complete receive a fixed penalty of 7,200 seconds, the cap on completed runs.
\end{itemize}
Each family or workload is run with five fixed seeds (11, 23, 37, 53, 71). Because seeds and workloads of the same software are not independent systems, we group cases by software engine (eight groups). For our primary grouping we also merge Spark and Hadoop into one ecosystem, since Spark runs on the Hadoop stack and the two share workload types, which leaves seven groups. Every family-level summary gives each group equal weight.

\paragraph{Native cohort.} To check that the recorded results are not an artifact of stored measurements, we also executed six systems natively on one machine (an Apple M3 Pro with 18\,GB of memory):
\begin{itemize}
  \item the cvc5 SMT solver \cite{barbosa2022cvc5} and the OR-Tools CP-SAT solver \cite{perron2023ortools}, both on N-queens instances;
  \item ripgrep \cite{gallant2016ripgrep}, running five fixed queries over the CPython source tree;
  \item hnswlib \cite{malkov2020hnsw}, building and querying an index over the Optdigits data, with at least 95\% recall required;
  \item Polars \cite{vink2020polars}, running queries over a public flights dataset;
  \item XGBoost \cite{chen2016xgboost}, training on Covertype, with at least 75\% validation accuracy required.
\end{itemize}
Each system has a finite grid of 64 configurations (128 for XGBoost). Every acquired configuration was checked for correctness or quality, and violations received a fixed penalty.

\begin{table}[t]
\centering
\caption{Recorded cohort. Every case has five seeds. Spark and Hadoop form one ecosystem in the primary analysis.}
\label{tab:cohort}
\small
\begin{tabular}{llllr}
\toprule
Engine & Source & Objective & Workloads & Cases \\
\midrule
BerkeleyDB (C) & DeepPerf / SPLConqueror & time (min.) & 1 & 5 \\
DUNE & DeepPerf / SPLConqueror & time (min.) & 1 & 5 \\
HIPAcc & DeepPerf / SPLConqueror & time (min.) & 1 & 5 \\
LLVM & DeepPerf / SPLConqueror & time (min.) & 1 & 5 \\
SaC & DeepPerf / SPLConqueror & time (min.) & 1 & 5 \\
OpenVPN & Performance Evolution & throughput (max.) & 1 & 5 \\
Spark & Tuneful & duration (min.) & 5 & 25 \\
Hadoop MapReduce & Scout & capped duration (min.) & 3 & 15 \\
\midrule
Total & & & & 70 \\
\bottomrule
\end{tabular}
\end{table}

\subsection{Paired-prefix protocol}

Every case starts from a saved prefix of ten labeled configurations.
\begin{itemize}
  \item For the six SPLConqueror and Performance Evolution families, the prefix consists of four configurations drawn in a seeded order followed by an EZR-style centroid search.
  \item For Spark and Hadoop, it consists of four seeded configurations followed by sequential three-nearest-neighbour (3NN) search.
\end{itemize}
Each continuation clones the prefix and acquires ten further configurations it has not seen before, for a total of 20 labels (B20). A continuation never sees another continuation's labels or any value it has not acquired. The candidate feature table is visible (the setting is transductive), but objective values are revealed only through a charged oracle. A continuation's result is the best value among its 20 labels.

The native cohort uses the same ten-label prefix but splits the remaining budget into seven search evaluations and three fresh measurements of the selected configuration, whose median is the reported result. All selected configurations, including those of older runs, were re-measured together in randomized blocks, so every comparison uses measurements from the same session.

\subsection{Classical continuations}

From each prefix we ran several cheap continuations:
\begin{itemize}
  \item \emph{sequential 3NN}, the reference, which continues from the prefix with nearest-neighbour search (for Spark and Hadoop this is the same optimizer that built the prefix);
  \item an \emph{adaptive neighbour} rule that searches around the current best configuration, and a \emph{fixed neighbour} rule that searches around the prefix's best;
  \item uniform \emph{random} search over unseen configurations;
  \item Gaussian-process expected improvement (\emph{GP-EI}) \cite{jones1998ego}, with a fixed mixed Mat\'ern-5/2 and Hamming kernel and targets normalized using acquired labels only;
  \item for Spark and Hadoop, \emph{random prototypes}: random proposals in the LLM's own proposal format, mapped to configurations by the same projection the LLM uses.
\end{itemize}
In the native cohort we also executed the acquisition functions of EZR version 0.9.4 without modification (the centroid and naive-Bayes variants) through a thin wrapper that injects the saved prefix \cite{menzies2026ezr}.

\subsection{LLM continuations}

\paragraph{Models.} We used three local models in 4-bit quantization (Q4\_K\_M), served with llama.cpp \cite{gerganov2023llamacpp}: SmolLM3-3B \cite{huggingface2025smollm3}, Qwen3-8B and Qwen3-14B \cite{qwen2025qwen3}. Thinking mode was off, the temperature was 0.7 and top-$p$ was 0.95. In the final experiment we added gpt-oss-120b \cite{openai2025gptoss}, the model used by SNAP2. We accessed it through OpenRouter \cite{openrouter2025} with the reasoning effort set to medium, the provider fixed for the whole run, and at most 16,000 output tokens for the one-shot interface and 8,000 for the loop.

\paragraph{One-shot interface.} The model receives the feature names and their allowed values, the ten measured configurations with their objective values, and the direction of optimization. It is asked for ten new configurations, written as fixed-width symbol strings. Local decoding is constrained by a grammar; hosted decoding is constrained by an equivalent JSON schema. Each proposal is projected to the nearest configuration not yet measured (Hamming distance for symbolic features; normalized distance plus category mismatch for Spark and Hadoop). In the native cohort, where only seven search evaluations remain, the first seven proposals are used and projection uses ordinal distance. A response that is malformed or truncated triggers the classical fallback, and the fallback's result is charged to the LLM arm.

\paragraph{Iterative loop.} To stay close to SNAP2 we also ran gpt-oss-120b in a loop of five rounds with two proposals per round. Each prompt contains:
\begin{itemize}
  \item the original task description;
  \item hard constraints on allowed symbols;
  \item every configuration measured so far, sorted from worst to best;
  \item a list of earlier proposals that collided with already-measured configurations;
  \item on a retry, the reason the previous answer was rejected.
\end{itemize}
Each round allows one retry, and a round that still fails falls back to the classical rule. We could not find a released SNAP2 implementation, so this loop is our adaptation of the design described in the paper, not a reproduction of its code.

\subsection{Measures}

For a case with reference result $y_{\mathrm{seq}}$ and continuation result $y$, the relative gain is $g = (y_{\mathrm{seq}} - y)/y_{\mathrm{seq}}$ when lower is better, with the sign reversed for throughput. Positive values favour the continuation. Because relative gains are asymmetric, we also report the mean log-ratio. We call a case \emph{useful} if $g > 1\%$.

The central measure is the \emph{LLM-specific win}: a case in which the LLM continuation beats sequential 3NN, random search, the adaptive neighbour rule and GP-EI each by more than 1\%. Its primary summary is the number of ecosystems, out of seven, that contain at least one LLM-specific win. We also report \emph{hindsight headroom}, the mean of $\max(0, g)$. It is an upper bound on what a perfect router could gain over the reference and is not a deployable policy.

For native systems we additionally require a gain above 10\%, different configurations, valid quality, stable repeated measurements (relative median absolute deviation of at most 5\%) and medians of at least 10\,ms. We call these \emph{robust} wins.

We do not report significance tests. Seven groups, one of which contains 40 of the 70 cases, cannot support population-level inference, and seeds are not independent systems.

\subsection{Controllers}

All controllers decide once, after the prefix, and use only information available at that point: 18 features describing the problem and the prefix trajectory. These include the number of candidates and options, spread and dispersion of the observed labels, progress between early and late prefix steps, plateau length, and a bootstrap estimate of how stable the incumbent is.

We trained four fixed benefit predictors on the paired gains: ridge regression on 7 or 18 features, a depth-2 regression tree, and RBF kernel ridge. Evaluation used nested leave-one-group-out cross-validation. Scaling, model fitting and the decision threshold were all chosen on inner folds that exclude the held-out group. We also report a fixed 80th-percentile threshold for each predictor, a bootstrap-uncertainty rule, and two rules adapted from prior work:
\begin{itemize}
  \item a BORA-inspired rule that calls the LLM when a GP is uncertain or the search has plateaued \cite{cisse2025bora};
  \item an LB-MCTS-inspired rule that calls the LLM with a probability derived from the surrogate's cross-validated rank reliability \cite{xu2026lbmcts}.
\end{itemize}
Both are checkpoint adaptations for a single decision, not the full iterative algorithms. We also transferred the frozen controllers, unchanged, to systems added later: WavPack, FFTW, Hadoop, Memcached and the six native systems.

\subsection{How the study evolved}
\label{sec:history}

The study went through many numbered iterations on the same systems. Along the way it changed prompts, output formats, datasets, budgets and controllers. We kept every protocol, raw model response, acquisition log and failure, and the replication package contains them all. Because the earlier analyses were carried out on data we had already seen, we treat them as exploratory.

The last two experiments were run differently. Before any data were collected, we froze the code, inputs and a decision rule by content hash:
\begin{itemize}
  \item a test of Qwen3-14B against a repeated Qwen3-8B draw;
  \item a test of gpt-oss-120b in the one-shot and loop interfaces.
\end{itemize}
Deviations during these runs were written down before any outcome was scored:
\begin{itemize}
  \item a failed server start that sent no requests;
  \item an output cap too small for a reasoning model, together with provider rate limits;
  \item wall-clock limits on individual stages.
\end{itemize}
The repeated Qwen3-8B draw hit its pre-declared memory cap and completed only 45 of its 70 cases. We report it only as a check on sampling variability.

\section{Results}

\subsection{\RQ{1}: Does the LLM continuation beat continued classical search?}
\label{sec:rq1}

It does not, on average. Table~\ref{tab:rq1} summarizes the recorded cohort. Every LLM arm has a negative mean gain relative to sequential 3NN, from $-3.1\%$ for gpt-oss-120b in the loop to $-5.2\%$ for Qwen3-8B, and the log-ratio agrees in sign everywhere. The cheap classical alternatives also lose to sequential 3NN on average, but by less. GP-EI is the strongest of them, with 22 useful cases against 11 cases more than 1\% worse.

\begin{table}[t]
\centering
\caption{\RQ{1}, recorded cohort (70 cases, seven ecosystems weighted equally). Gains are relative to continued sequential 3NN from the same prefix. ``Better'' and ``worse'' count cases above and below a 1\% margin. Headroom is the mean of $\max(0,g)$.}
\label{tab:rq1}
\small
\begin{tabular}{lrrrrr}
\toprule
Continuation & Mean gain & Log-ratio & Better & Worse & Headroom \\
\midrule
Random search & $-3.49\%$ & $-2.97\%$ & 11 & 22 & 0.32\% \\
Adaptive neighbour & $-2.29\%$ & $-1.90\%$ & 10 & 17 & 0.48\% \\
GP-EI & $-2.19\%$ & $-1.66\%$ & 22 & 11 & 1.28\% \\
\midrule
SmolLM3-3B & $-4.96\%$ & $-4.56\%$ & 10 & 24 & 0.22\% \\
Qwen3-8B & $-5.24\%$ & $-4.79\%$ & 5 & 29 & 0.18\% \\
Qwen3-14B & $-4.00\%$ & $-3.60\%$ & 9 & 23 & 0.45\% \\
gpt-oss-120b, one-shot, draw 1 & $-3.73\%$ & $-3.48\%$ & 15 & 14 & 0.53\% \\
gpt-oss-120b, one-shot, draw 2 & $-3.16\%$ & $-2.69\%$ & 19 & 15 & 0.76\% \\
gpt-oss-120b, iterative loop & $-3.14\%$ & $-2.49\%$ & 14 & 14 & 0.87\% \\
\bottomrule
\end{tabular}
\end{table}

The averages hide large differences between ecosystems (Table~\ref{tab:eco}). DUNE is hard for every continuation: from these prefixes nothing comes within 10\% of sequential 3NN. OpenVPN punishes the LLMs in particular, with losses of 5.9--9.0\% for four of the six LLM arms against 1--2\% for the classical alternatives. Spark and Hadoop form the only ecosystem where several LLM arms gain more than 2\% on average, and there GP-EI gains too. Elsewhere, LLM gains, where they occur at all, stay below about 1\%.

\begin{table}[t]
\centering
\caption{Mean gain relative to sequential 3NN by ecosystem (recorded cohort). The Spark/Hadoop row averages 40 cases; every other row averages 5.}
\label{tab:eco}
\small
\setlength{\tabcolsep}{4pt}
\begin{tabular}{lrrrrrrrrr}
\toprule
 & \multicolumn{3}{c}{Classical} & \multicolumn{3}{c}{Local LLM} & \multicolumn{3}{c}{gpt-oss-120b} \\
\cmidrule(lr){2-4}\cmidrule(lr){5-7}\cmidrule(lr){8-10}
Ecosystem & Random & Adapt. & GP-EI & 3B & 8B & 14B & Draw 1 & Draw 2 & Loop \\
\midrule
BerkeleyDB & $-1.04$ & $+0.71$ & $+4.51$ & $-0.33$ & $+0.06$ & $-0.87$ & $-0.34$ & $+0.92$ & $+0.21$ \\
DUNE & $-18.14$ & $-11.15$ & $-18.06$ & $-20.89$ & $-21.24$ & $-18.61$ & $-18.01$ & $-18.38$ & $-20.63$ \\
HIPAcc & $-0.64$ & $-0.71$ & $-1.06$ & $-2.60$ & $-2.91$ & $-1.37$ & $+1.03$ & $-0.89$ & $-1.61$ \\
LLVM & $-2.20$ & $+0.14$ & $-1.61$ & $-1.97$ & $-2.16$ & $-1.00$ & $-0.01$ & $-1.07$ & $-1.62$ \\
OpenVPN & $-0.97$ & $-2.23$ & $-2.19$ & $-7.08$ & $-6.88$ & $-5.87$ & $-9.00$ & $-4.59$ & $-2.38$ \\
SaC & $-0.27$ & $-0.40$ & $-0.14$ & $-0.40$ & $-0.40$ & $-0.27$ & $-0.14$ & $-0.40$ & $+0.13$ \\
Spark/Hadoop & $-1.15$ & $-2.41$ & $+3.18$ & $-1.42$ & $-3.16$ & $+0.01$ & $+0.37$ & $+2.29$ & $+3.93$ \\
\bottomrule
\end{tabular}
\end{table}

The native cohort tells the same story with SmolLM3-3B and Qwen3-8B (Table~\ref{tab:native}). On OR-Tools and XGBoost the LLM continuations lose by 20--26\%. On cvc5, ripgrep and hnswlib the differences are within a few percent. In 17 of the 30 case pairs the LLM chose exactly the configuration that sequential 3NN chose. No continuation of any kind, classical or LLM, produced a robust win over sequential 3NN on these 64- and 128-configuration spaces. After ten of 64 evaluations there was little left to find.

\begin{table}[t]
\centering
\caption{Native cohort: mean gain relative to sequential 3NN (five seeds each; medians of three fresh measurements taken in the same session). No arm, classical or LLM, has a robust win on any system.}
\label{tab:native}
\small
\begin{tabular}{lrrr}
\toprule
System & SmolLM3-3B & Qwen3-8B & Robust wins (any arm) \\
\midrule
cvc5 (N-queens) & $-0.11\%$ & $+0.17\%$ & 0 \\
OR-Tools CP-SAT (N-queens) & $-20.67\%$ & $-20.22\%$ & 0 \\
ripgrep & $-2.17\%$ & $-0.32\%$ & 0 \\
hnswlib & $-0.15\%$ & $-0.24\%$ & 0 \\
Polars & $-5.07\%$ & $-4.83\%$ & 0 \\
XGBoost & $-26.11\%$ & $-19.53\%$ & 0 \\
\bottomrule
\end{tabular}
\end{table}

\subsection{\RQ{2}: Are the LLM wins specific to the LLM?}
\label{sec:rq2}

For the local models, mostly not. Table~\ref{tab:rq1} already hints at this: random search has as many useful cases as SmolLM3-3B and more than Qwen3-8B. The case-level picture is starker. For all 10 of SmolLM3-3B's useful cases, and for 4 of Qwen3-8B's 5, at least one of random search, the adaptive neighbour rule or GP-EI was also more than 1\% better than sequential 3NN. The LLM was mostly winning where the incumbent optimizer had stalled, and where almost any change of strategy helped.

The random-prototype control makes the point most directly. In Spark and Hadoop, random proposals passed through exactly the same projection as the LLM's proposals beat sequential 3NN on average (+1.17\% on Spark, +4.42\% on Hadoop). SmolLM3-3B's own proposals did worse there ($-0.11\%$ and $-3.60\%$).

Table~\ref{tab:rq2} counts LLM-specific wins. SmolLM3-3B and Qwen3-8B had none, and Qwen3-14B had two, both on Spark. Random proposals passed through the same projection also reached one of those two results. The same holds in the native cohort, where no model had a case that beat every classical alternative by more than 1\%. The LLMs' hindsight headroom there (0.36\% and 0.40\%) was barely above random search's (0.30\%) and less than half that of EZR's naive-Bayes acquisition (0.90\%).

\begin{table}[t]
\centering
\caption{\RQ{2}--\RQ{3}: LLM-specific wins (recorded cohort). A win beats sequential 3NN, random search, the adaptive neighbour rule and GP-EI each by more than 1\%. The last three columns count cases in which the LLM beats each single alternative by more than 1\%.}
\label{tab:rq2}
\small
\begin{tabular}{lrrrrr}
\toprule
 & \multicolumn{2}{c}{LLM-specific wins} & \multicolumn{3}{c}{Beats one alternative} \\
\cmidrule(lr){2-3}\cmidrule(lr){4-6}
LLM arm & Ecosystems (of 7) & Cases & Random & Adapt. & GP-EI \\
\midrule
SmolLM3-3B & 0 & 0 & 12 & 11 & 5 \\
Qwen3-8B & 0 & 0 & 4 & 6 & 4 \\
Qwen3-14B & 1 & 2 & 11 & 11 & 9 \\
gpt-oss-120b, one-shot, draw 1 & 2 & 3 & 22 & 15 & 7 \\
gpt-oss-120b, one-shot, draw 2 & 1 & 1 & 26 & 18 & 10 \\
gpt-oss-120b, iterative loop & 1 & 2 & 21 & 18 & 9 \\
\bottomrule
\end{tabular}
\end{table}

\subsection{\RQ{3}: Do stronger models or iterative interaction help?}
\label{sec:rq3}

Somewhat. Moving from Qwen3-8B to Qwen3-14B, with prompts and seeds unchanged and rendered prompts identical byte for byte, reduced the mean loss from 5.24\% to 4.00\%. Qwen3-14B was more than 1\% better than the earlier 8B draw in 16 cases and worse in 5. gpt-oss-120b went further. In its first one-shot draw it was more than 1\% better than Qwen3-8B in 27 cases and worse in 5. It was also the only model with LLM-specific wins in more than one ecosystem.

Before collecting the gpt-oss-120b data we fixed a decision rule: LLM-specific wins in at least two of the seven ecosystems in either the first one-shot draw or the loop would count as evidence of LLM-specific headroom at this model scale. The first draw met the rule exactly (Table~\ref{tab:rq2}), so formally the answer is yes. The result is weaker than that sounds:
\begin{itemize}
  \item The two HIPAcc cases that supplied the second ecosystem won by only 1.9\% and 1.7\% and did not recur. In the second draw they scored $-5.6\%$ and 0.0\%; in the loop, $-1.7\%$ and $-2.0\%$.
  \item The second draw and the loop each reached only one ecosystem.
  \item When we raised the margin to 2\% after seeing the data, or required a win to appear in at least two of the three runs, every run fell to one ecosystem, Spark/Hadoop.
\end{itemize}

One case is solid. On the Bayes workload of Spark with seed 11, gpt-oss-120b beat every cheap alternative in all three runs:
\begin{itemize}
  \item by 8.3\% over sequential 3NN in both one-shot draws;
  \item by 31.5\% over sequential 3NN (and 27\% over GP-EI) in the loop.
\end{itemize}
None of the local models found it. In a second Spark case (TeraSort, seed 23) the loop reached a result 24\% better than sequential 3NN, matching what Qwen3-14B had found.

The iterative loop was the best way we found to use gpt-oss-120b. It had the best mean gain and the largest headroom (Table~\ref{tab:rq1}), and against the first one-shot draw it was more than 1\% better in 19 cases and worse in 10. This agrees with SNAP2's design choice. It contrasts with an earlier exploratory experiment with Qwen3-8B on LLVM and SaC, in which five rounds of measured feedback and a control with the feedback masked ended with the same incumbent in all ten cases. Feedback seems to help only when the model can use it.

\subsection{\RQ{4}: Can a controller predict useful escalations?}
\label{sec:rq4}

Not on this evidence, and we think it could not have. Table~\ref{tab:rq4} shows the controllers on the recorded cohort for the two models that had enough paired data to train on.

\begin{table}[t]
\centering
\caption{\RQ{4}: controllers under leave-one-ecosystem-out evaluation (70 cases per model). The gain is relative to sequential 3NN and is zero for any case in which the controller does not call the LLM. The oracle picks the better branch after seeing both outcomes and is not deployable.}
\label{tab:rq4}
\small
\begin{tabular}{lrrrr}
\toprule
 & \multicolumn{2}{c}{Calls (of 70)} & \multicolumn{2}{c}{Mean gain} \\
\cmidrule(lr){2-3}\cmidrule(lr){4-5}
Policy & SmolLM3 & Qwen3-8B & SmolLM3 & Qwen3-8B \\
\midrule
Never call & 0 & 0 & $0.00\%$ & $0.00\%$ \\
Always call & 70 & 70 & $-4.96\%$ & $-5.24\%$ \\
Benefit predictor (inner-CV selection) & 0 & 0 & $0.00\%$ & $0.00\%$ \\
Ridge, 18 features, 80th percentile & 9 & 9 & $-0.10\%$ & $-0.10\%$ \\
Regression tree, 80th percentile & 4 & 8 & $-2.97\%$ & $-0.45\%$ \\
Bootstrap uncertainty, 80th percentile & 27 & 27 & $-3.12\%$ & $-3.16\%$ \\
BORA-inspired checkpoint rule & 14 & 14 & $-3.10\%$ & $-3.25\%$ \\
Rank-reliability rule (expected calls) & 23.95 & 23.95 & $-1.44\%$ & $-1.49\%$ \\
\midrule
Hindsight oracle & -- & -- & $+0.22\%$ & $+0.18\%$ \\
\bottomrule
\end{tabular}
\end{table}

When the threshold was chosen on inner folds, the learned benefit predictor never called the LLM. The same happened when the frozen controllers were transferred to systems added later: WavPack, FFTW, Hadoop, Memcached and the six native systems. Every rule that did call the LLM lost, and the uncertainty, BORA-inspired and rank-reliability rules also did worse than randomly chosen calls at the same rate. In the native cohort the two adapted rules made between 1 and 7 calls per ten cases. Their outcomes relative to never calling ranged from a gain of 0.15\% (from a single call) to a loss of 7.5\%.

We initially read these results as failures of the controllers. The structure of the data suggests otherwise. For SmolLM3-3B, all ten useful cases fall in the Spark/Hadoop ecosystem:
\begin{itemize}
  \item holding that ecosystem out leaves the training folds with no positive examples;
  \item holding out any other ecosystem leaves the test fold with nothing to find.
\end{itemize}
None of the seven leave-one-ecosystem-out folds could reward a working predictor. Qwen3-8B had two such folds out of seven. For LLM-specific wins the picture is even sparser: at most two ecosystems contain one for any model, including gpt-oss-120b.

Under these conditions a threshold rule tuned for gain on the training groups must end up choosing never to call, because calling loses on average everywhere it can be tested. The hindsight oracle's headroom, below 0.25\% for both local models, bounds what any controller could have gained. That bound is smaller than the hindsight headroom of simply switching to GP-EI (1.28\%).

\section{Discussion}

\subsection{What this means for practitioners}

With a budget of around 20 evaluations and after a classical warm-up, calling an LLM by default is not a good bet. This held for models from 3B to about 120B parameters, both local and hosted. If the incumbent optimizer has stalled, switching to a different cheap strategy (GP-EI in our data) captures most of what an LLM would have gained, at a small fraction of the cost. When LLM escalation paid off in our data, it did so in a handful of cases on one family of big-data systems. A stronger model used iteratively seems to be needed for those gains, and even then they are unpredictable from the prefix.

\subsection{What this means for researchers}

\paragraph{Choose the right counterfactual.} Improvement over the prefix, or over the incumbent optimizer's own continuation, overstates the value of an LLM. In our data, measuring against the best cheap alternative turned ``the LLM sometimes helps'' into ``the 3B and 8B models never helped in a way that other cheap strategies could not match, and the 14B model did so twice''. We recommend reporting at least random search, a neighbourhood rule and a Bayesian-optimization baseline from the same prefix and budget. The budget-matched study by Rodrigues et al.\ reaches a similar conclusion for hyperparameter optimization \cite{rodrigues2026worth}.

\paragraph{Check that a router can be tested before building one.} Leave-one-group-out evaluation of a router is only informative when several groups contain positive cases. Counting the positive cases per group should come before any router is trained. A router that learns never to call the LLM should then be read in light of that count: here, never calling was correct.

\paragraph{Averages hide rare wins.} The mean gain of gpt-oss-120b was negative in every run, yet it produced a genuine win of up to 31\% that no other method found. A benefit that appears rarely and in only a few places matters for practice, but it is hard to learn or to predict, and it makes sample size a pressing concern for evaluation.

\paragraph{Interaction design matters as much as model size.} The loop outperformed one-shot use of the same model. The earlier null result with Qwen3-8B suggests that the model must be strong enough to use feedback before the loop pays off.

\subsection{Cost}

Classical continuations cost milliseconds of computation on recorded data. A local LLM escalation cost roughly 1--3\,s per request for SmolLM3-3B, 3--15\,s for Qwen3-8B and about 26\,s for Qwen3-14B on our laptop; Qwen3-14B used 11.3\,GB of memory at its peak. A hosted gpt-oss-120b call took about 26--30\,s and cost about \$0.0035. It typically generated 4,500--7,800 output tokens, most of them reasoning. The loop cost about \$0.013 per case. The whole hosted experiment, 513 successful calls, cost \$1.38. These prices are low in absolute terms, but they buy no average benefit, and the wall-clock time per case is far larger than that of the classical alternatives.

\section{Threats to Validity}

\paragraph{Construct validity.}
\begin{itemize}
  \item Relative gains are asymmetric, so we report log-ratios, which agree in sign.
  \item The 1\% margin is a choice. Post-hoc checks at 2\% and 5\% make the gpt-oss-120b result weaker, not stronger.
  \item Recorded measurements are single runs with no noise estimate. The native cohort, with repeated measurements, supports the same conclusions.
  \item For Hadoop, failed runs receive a fixed penalty rather than a measured runtime.
  \item The hosted model's weights cannot be verified by hash. We recorded the provider, model identifier, response identifiers and token usage for every call.
  \item SNAP2 does not report its reasoning setting, so we chose medium effort in advance.
\end{itemize}

\paragraph{Internal validity.}
\begin{itemize}
  \item The study's long history on the same systems creates the risk of a garden of forking paths \cite{gelman2014forking}. Most of those forks tried to make the LLM arm work better, which works against our negative findings rather than for them.
  \item The last two experiments were frozen before collection, and every deviation was recorded before outcomes were scored.
  \item The repeated Qwen3-8B draw hit its memory cap and is incomplete.
  \item Hosted inference is not deterministic, even with fixed seeds.
\end{itemize}

\paragraph{External validity.}
\begin{itemize}
  \item Seven ecosystems is a small sample, and one of them supplies 40 of the 70 cases.
  \item The native systems have small configuration spaces that left no robust headroom for any method.
  \item We used one budget (B20) and one hand-off point (ten evaluations).
  \item All local runs used one laptop, and all hosted runs one provider.
  \item All tasks use public data that may appear in the models' pretraining data.
  \item Our results do not rule out benefit at other budgets or with other models or interfaces, and they do not refute SNAP2's results on its own tasks.
\end{itemize}

\paragraph{Conclusion validity.}
\begin{itemize}
  \item We report no significance tests, and we weight ecosystems equally.
  \item Seeds are repeated measurements of the same system, not independent observations.
  \item Instability across repeated runs is itself a finding \cite{rayegan2026instability}: two draws of the same model on the same prompts disagreed by more than 1\% in 26 of 70 cases.
\end{itemize}

\section{Conclusion}

We asked whether a cheap controller could tell, after a short classical search, when an LLM continuation would be worth calling. We found that there was little to predict. Local models from 3B to 8B parameters produced no gain that a cheap classical alternative could not match, and the 14B model managed it twice. SNAP2's own model produced a few genuine gains, one of them large and repeatable, but concentrated in a single ecosystem and outweighed on average by losses. Under those conditions no controller can be shown to work across systems, and abstaining is the correct policy.

We think the more durable lessons are methodological. Escalation benefit should be measured against the best cheap alternative, and routing studies should first establish that benefit exists in several independent groups. A cohort of fresh systems with large configuration spaces, in which LLM-specific benefit appears in many groups, would be needed to test a controller properly. Assembling one is the obvious next step.

\section*{Data Availability}

Protocols, frozen input hashes, raw model responses and token usage, acquisition logs, analysis and verification scripts, and integrity manifests for all experiments are available in the replication package [URL to be added]. The package contains no model weights or API credentials.

\section*{Acknowledgements}

[To be added.]

\bibliographystyle{plainnat}
\bibliography{references}

\end{document}
